\chapter{Phân loại trạng thái\\{\normalsize\textit{Classification of States}}}
\label{ch:classification}

\textbf{Câu chuyện:} Hãy tưởng tượng bạn đi du lịch giữa các thành phố.
Có những thành phố bạn chắc chắn sẽ quay lại (quán café quen --- ``hồi quy''),
và có những thành phố bạn chỉ ghé qua một lần rồi không bao giờ quay lại (``tạm thời'').
Chương này giúp phân loại: trạng thái nào ``quen'' (hồi quy), trạng thái nào ``lạ'' (tạm thời)?

% ============================================================================
\section{Trạng thái liên thông (Communicating States)}
\label{sec:communicating-states}

\textbf{Câu chuyện:} Hãy tưởng tượng bạn đang ở thành phố $i$. Câu hỏi đặt ra:
``Liệu bạn có thể đi đến thành phố $j$ được không?''
Nếu có một con đường (dù phải đi qua nhiều thành phố trung gian),
tức là tồn tại số bước $n$ sao cho $P_{i,j}^{(n)} > 0$,
thì ta nói $j$ \textbf{có thể đạt đến} từ $i$.

Ví dụ đời thường: Từ Hà Nội bạn có thể bay đến Tokyo (qua Seoul chuyển tiếp)
--- vậy Tokyo ``có thể đạt đến'' từ Hà Nội.

\begin{dinhnghia}[(Trạng thái có thể đạt đến -- \en{Accessible state})]
\label{def:accessible}
Cho chuỗi Markov $(X_n)_{n \in \N}$ với không gian trạng thái $S$ và ma trận chuyển $P$.
Trạng thái $j$ được gọi là \term{có thể đạt đến} (\en{accessible}) từ trạng thái $i$,
ký hiệu $i \to j$, nếu tồn tại $n \geq 0$ sao cho:
\[
  P_{i,j}^{(n)} := \PP(X_n = j \mid X_0 = i) > 0.
\]
Nói cách khác, xuất phát từ $i$, chuỗi Markov có xác suất dương để đến được $j$
sau một số bước hữu hạn nào đó.
\end{dinhnghia}

\begin{dinhnghia}[(Trạng thái liên thông -- \en{Communicating states})]
\label{def:communicating}
Hai trạng thái $i$ và $j$ được gọi là \term{liên thông} (\en{communicate}),
ký hiệu $i \leftrightarrow j$, nếu $i \to j$ và $j \to i$, tức là:
\[
  \exists\, a \geq 0: \; P_{i,j}^{(a)} > 0
  \qquad \text{và} \qquad
  \exists\, b \geq 0: \; P_{j,i}^{(b)} > 0.
\]
\end{dinhnghia}

Ví dụ: Hà Nội và Tokyo ``liên thông'' vì bạn có thể bay đi lẫn bay về.
Nhưng nếu đảo $X$ chỉ có chuyến bay đến mà không có chuyến bay đi,
thì $X$ ``có thể đạt đến'' nhưng không ``liên thông'' với các thành phố khác.

Quan hệ ``liên thông'' $\leftrightarrow$ là một \term{quan hệ tương đương} (\en{equivalence relation})
trên $S$, vì nó thỏa mãn ba tính chất:

\begin{itemize}[nosep]
  \item \textbf{Phản xạ} (\en{Reflexive}): $i \leftrightarrow i$ vì $P_{i,i}^{(0)} = 1 > 0$.
  \item \textbf{Đối xứng} (\en{Symmetric}): $i \leftrightarrow j \Rightarrow j \leftrightarrow i$ (theo định nghĩa).
  \item \textbf{Bắc cầu} (\en{Transitive}): Nếu $i \leftrightarrow j$ và $j \leftrightarrow k$ thì $i \leftrightarrow k$.
        Thật vậy, nếu $P_{i,j}^{(a)} > 0$ và $P_{j,k}^{(b)} > 0$ thì theo phương trình Chapman--Kolmogorov:
        \[
          P_{i,k}^{(a+b)} = \sum_{l \in S} P_{i,l}^{(a)} P_{l,k}^{(b)}
          \geq P_{i,j}^{(a)} P_{j,k}^{(b)} > 0.
        \]
\end{itemize}

Vì $\leftrightarrow$ là quan hệ tương đương, nó \term{phân hoạch} (\en{partition}) không gian trạng thái
$S$ thành các \term{lớp tương đương} (\en{equivalence classes}), gọi là \term{lớp liên thông}
(\en{communicating classes}):
\begin{equation}\label{eq:partition}
  S = A_1 \cup A_2 \cup \cdots \cup A_M,
  \qquad A_k \cap A_l = \emptyset \;\text{với}\; k \neq l.
\end{equation}

\begin{tomtat}
Quan hệ $\leftrightarrow$ phân chia không gian trạng thái $S$ thành các lớp liên thông rời nhau.
Bên trong mỗi lớp, mọi cặp trạng thái đều có thể đạt đến nhau.
Giữa các lớp khác nhau, việc chuyển đổi có thể chỉ xảy ra theo một chiều hoặc không xảy ra.

Ví dụ: trong mạng lưới giao thông, các thành phố trong cùng khu vực đều liên thông.
Nhưng ``vùng hấp thụ'' (ví dụ: hố đen) chỉ có đường vào mà không có đường ra.
\end{tomtat}

% ============================================================================
\section{Lớp liên thông (Communicating Class)}
\label{sec:communicating-class}

\begin{dinhnghia}[(Chuỗi bất khả quy -- \en{Irreducible chain})]
\label{def:irreducible}
Một chuỗi Markov được gọi là \term{bất khả quy} (\en{irreducible}) nếu toàn bộ không gian
trạng thái $S$ chỉ gồm duy nhất một lớp liên thông, tức là mọi cặp trạng thái $i, j \in S$
đều liên thông: $i \leftrightarrow j$.

Ngược lại, nếu $S$ có ít nhất hai lớp liên thông thì chuỗi được gọi là
\term{khả quy} (\en{reducible}).
\end{dinhnghia}

\textbf{Ví dụ trực quan:}
\begin{itemize}[nosep]
  \item \textbf{Bất khả quy:} Bạn có thể đi từ bất kỳ thành phố nào đến bất kỳ thành phố nào (dù qua trung chuyển).
  \item \textbf{Khả quy:} Có ``vùng cấm'' --- một khi vào rồi, không ra được.
\end{itemize}

\begin{example}[Chuỗi khả quy -- \en{Reducible chain}]
\label{ex:reducible-chain}
Xét chuỗi Markov trên $S = \{0, 1, 2, 3\}$ với ma trận chuyển:
\begin{equation}\label{eq:reducible-matrix}
  P = \begin{bmatrix}
    0   & 0.2 & 0.8 & 0   \\
    0.3 & 0.1 & 0   & 0.6 \\
    0.5 & 0   & 0   & 0.5 \\
    0   & 0   & 0   & 1
  \end{bmatrix}.
\end{equation}

\textbf{Phân tích:}
\begin{itemize}[nosep]
  \item Từ trạng thái $0$: có thể đi đến $1$ (qua $P_{0,1} = 0.2$) và đến $2$ (qua $P_{0,2} = 0.8$).
  \item Từ trạng thái $1$: có thể đi đến $0$ (qua $P_{1,0} = 0.3$) và đến $3$ (qua $P_{1,3} = 0.6$).
  \item Từ trạng thái $2$: có thể đi đến $0$ (qua $P_{2,0} = 0.5$) và đến $3$ (qua $P_{2,3} = 0.5$).
  \item Từ trạng thái $3$: $P_{3,3} = 1$, trạng thái $3$ là \term{hấp thụ} (\en{absorbing}).
\end{itemize}

Các trạng thái $0, 1, 2$ liên thông với nhau, nhưng từ $3$ không thể quay lại $0, 1, 2$.
Do đó chuỗi có hai lớp liên thông:
\[
  A_1 = \{0, 1, 2\}, \qquad A_2 = \{3\}.
\]
Chuỗi này là \term{khả quy} (\en{reducible}).
\end{example}

\begin{tomtat}
Chuỗi \term{bất khả quy}: chỉ có một lớp liên thông, mọi trạng thái đều đạt đến được từ bất kỳ
trạng thái nào khác.

Chuỗi \term{khả quy}: có nhiều lớp liên thông; một khi vào một lớp ``đóng'' (\en{closed class}),
chuỗi không thể thoát ra.
\end{tomtat}

% ============================================================================
\section{Trạng thái hồi quy (Recurrent States)}
\label{sec:recurrent-states}

\textbf{Câu chuyện:} Quán café quen là ``hồi quy'' --- dù bạn đi đâu, chắc chắn một ngày sẽ quay lại.
Quán ghé qua lần đầu ở thành phố lạ là ``tạm thời'' --- có thể bạn không bao giờ quay lại nữa.

Nhớ lại từ Chương~\ref{ch:first-step}, \term{thời gian quay lại lần đầu} (\en{first return time}) trạng thái $i$
được định nghĩa là:
\[
  \tau_i := \inf\{n \geq 1 : X_n = i\},
\]
và \term{xác suất quay lại} (\en{return probability}):
\[
  p_{ii} := \PP(\tau_i < \infty \mid X_0 = i).
\]

\begin{dinhnghia}[(Trạng thái hồi quy -- \en{Recurrent state})]
\label{def:recurrent}
Trạng thái $i \in S$ được gọi là \term{hồi quy} (\en{recurrent}) nếu:
\begin{equation}\label{eq:recurrent-def}
  p_{ii} = \PP(\tau_i < \infty \mid X_0 = i) = 1.
\end{equation}
Tức là, xuất phát từ $i$, chuỗi Markov \emph{chắc chắn} sẽ quay lại $i$.
\end{dinhnghia}

Tính hồi quy có thể được đặc trưng bằng các điều kiện tương đương sau.

\begin{dinhly}[(Đặc trưng của trạng thái hồi quy)]
\label{thm:recurrent-equiv}
Các mệnh đề sau là tương đương:
\begin{enumerate}[label=(\roman*)]
  \item $i$ là trạng thái hồi quy, tức $p_{ii} = 1$.
  \item \term{Số lần quay lại trung bình} (\en{expected number of returns}) là vô hạn:
        \begin{equation}\label{eq:recurrent-Ri}
          \E[R_i \mid X_0 = i] = +\infty,
        \end{equation}
        trong đó $R_i = \sum_{n=1}^{\infty} \ind_{\{X_n = i\}}$ là tổng số lần chuỗi thăm $i$.
  \item $\PP(R_i = +\infty \mid X_0 = i) = 1$, tức chuỗi quay lại $i$ \emph{vô hạn lần} với xác suất~$1$.
  \item Chuỗi các xác suất quay lại phân kỳ:
        \begin{equation}\label{eq:recurrent-series}
          \sum_{n=0}^{\infty} P_{i,i}^{(n)} = +\infty.
        \end{equation}
\end{enumerate}
\end{dinhly}

\begin{congthuc}
\textbf{Tiêu chuẩn hồi quy (bằng chuỗi):}
\[
  i \text{ hồi quy} \quad\Longleftrightarrow\quad
  \sum_{n=0}^{\infty} P_{i,i}^{(n)} = +\infty.
\]
Trực quan: nếu ``tổng xác suất quay lại'' phân kỳ, thì chuỗi quay lại vô hạn lần.
\end{congthuc}

% ============================================================================
\section{Trạng thái tạm thời (Transient States)}
\label{sec:transient-states}

\begin{dinhnghia}[(Trạng thái tạm thời -- \en{Transient state})]
\label{def:transient}
Trạng thái $i \in S$ được gọi là \term{tạm thời} (\en{transient}) nếu:
\begin{equation}\label{eq:transient-def}
  \PP(R_i = +\infty \mid X_0 = i) < 1,
\end{equation}
hay tương đương, $p_{ii} < 1$.
\end{dinhnghia}

Đối với trạng thái tạm thời:
\begin{itemize}[nosep]
  \item $\E[R_i \mid X_0 = i] = \dfrac{p_{ii}}{1 - p_{ii}} < +\infty$: số lần quay lại trung bình là hữu hạn.
  \item $\displaystyle\sum_{n=0}^{\infty} P_{i,i}^{(n)} < +\infty$: chuỗi các xác suất quay lại hội tụ.
  \item Với xác suất $1 - p_{ii} > 0$, chuỗi sẽ \emph{rời khỏi} $i$ và không bao giờ quay lại.
\end{itemize}

\begin{tomtat}
\begin{center}
\renewcommand{\arraystretch}{1.4}
\begin{tabular}{lcc}
\hline
\textbf{Tính chất} & \textbf{Hồi quy} (\en{Recurrent}) & \textbf{Tạm thời} (\en{Transient}) \\
\hline
$p_{ii}$ & $= 1$ & $< 1$ \\
$\E[R_i \mid X_0 = i]$ & $= +\infty$ & $< +\infty$ \\
$\sum_n P_{i,i}^{(n)}$ & $= +\infty$ & $< +\infty$ \\
Số lần thăm $i$ & Vô hạn lần (h.c.c.) & Hữu hạn lần (h.c.c.) \\
Ví dụ đời thường & Quán café quen & Quán ghé qua khi du lịch \\
\hline
\end{tabular}
\end{center}
\end{tomtat}

% --- Corollary: recurrence is a class property ---

\begin{corollary}[(Hồi quy là tính chất lớp -- \en{Recurrence is a class property})]
\label{cor:recurrence-class}
Nếu trạng thái $i$ là hồi quy và $j$ liên thông với $i$ (tức $i \leftrightarrow j$),
thì $j$ cũng là hồi quy.
\end{corollary}

\begin{proof}
Vì $i \leftrightarrow j$, tồn tại $a, b \geq 0$ sao cho $P_{j,i}^{(a)} > 0$ và $P_{i,j}^{(b)} > 0$.
Theo phương trình Chapman--Kolmogorov, với mọi $n \geq 0$:
\[
  P_{j,j}^{(a+n+b)} \geq P_{j,i}^{(a)} \, P_{i,i}^{(n)} \, P_{i,j}^{(b)}.
\]
Do đó:
\[
  \sum_{m=0}^{\infty} P_{j,j}^{(m)}
  \geq P_{j,i}^{(a)} \, P_{i,j}^{(b)} \sum_{n=0}^{\infty} P_{i,i}^{(n)}
  = +\infty,
\]
vì $i$ là hồi quy nên $\sum_n P_{i,i}^{(n)} = +\infty$, và $P_{j,i}^{(a)} P_{i,j}^{(b)} > 0$.
Vậy $j$ cũng là hồi quy.
\end{proof}

\textbf{Hệ quả:} Tạm thời cũng là tính chất lớp. Nếu $i$ là tạm thời và $i \leftrightarrow j$,
thì $j$ cũng là tạm thời.

% --- Example: Simple random walk ---

\begin{example}[Bước đi ngẫu nhiên đơn giản trên $\Z$ -- \en{Simple random walk on $\Z$}]
\label{ex:random-walk-recurrence}
Xét bước đi ngẫu nhiên $(S_n)_{n \geq 0}$ trên $\Z$ với $S_0 = 0$,
$\PP(X_k = +1) = p$ và $\PP(X_k = -1) = q = 1-p$.

Xác suất trở về gốc tại thời điểm chẵn $n = 2k$:
\[
  P_{0,0}^{(2k)} = \binom{2k}{k} p^k q^k, \qquad P_{0,0}^{(2k+1)} = 0.
\]

\textbf{Trường hợp $p = q = 1/2$:} $4pq = 1$, suy ra $P_{0,0}^{(2k)} \sim \frac{1}{\sqrt{\pi k}}$ và:
\[
  \sum_{k=0}^{\infty} P_{0,0}^{(2k)} = +\infty
  \qquad \Longrightarrow \qquad
  \text{trạng thái $0$ là \textbf{hồi quy}.}
\]

\textbf{Trường hợp $p \neq q$:} $4pq < 1$, suy ra chuỗi hội tụ:
\[
  \sum_{k=0}^{\infty} P_{0,0}^{(2k)} < +\infty
  \qquad \Longrightarrow \qquad
  \text{trạng thái $0$ là \textbf{tạm thời}.}
\]

Trực quan: với $p \neq q$, bước đi ngẫu nhiên có ``xu hướng'' đi về một phía,
nên càng lúc càng xa gốc và khó quay lại.
\end{example}

\begin{congthuc}
\textbf{Bước đi ngẫu nhiên đơn giản trên $\Z$:}
\[
  \text{Hồi quy} \;\Longleftrightarrow\; p = q = \frac{1}{2}.
\]
\end{congthuc}

% --- Polya's theorem ---
\begin{dinhly}[(Định lý Polya về hồi quy -- \en{Polya's recurrence theorem})]
\label{thm:polya}
Bước đi ngẫu nhiên đơn giản đối xứng ($p = 1/2$ mỗi chiều) trên mạng lưới $\Z^d$:
\begin{itemize}[nosep]
  \item $d = 1$: \textbf{hồi quy} (chắc chắn quay lại gốc).
  \item $d = 2$: \textbf{hồi quy} (``A drunk man will find his way home'' --- Polya).
  \item $d \geq 3$: \textbf{tạm thời} (``A drunk bird will NOT find its way home'').
\end{itemize}
\end{dinhly}

Trực quan: trong không gian 1D và 2D, ``không có đủ chỗ để trốn'' nên bước đi phải quay lại.
Trong 3D trở lên, không gian quá rộng, bước đi ``lạc'' mãi mãi với xác suất dương.

% --- Theorem 4: Finite chain has at least one recurrent state ---

\begin{dinhly}[(Tồn tại trạng thái hồi quy trong chuỗi hữu hạn)]
\label{thm:finite-recurrent}
Mọi chuỗi Markov với không gian trạng thái hữu hạn $S$ có ít nhất một trạng thái hồi quy.
\end{dinhly}

\begin{proof}
Giả sử bằng phản chứng rằng \emph{tất cả} trạng thái đều tạm thời.
Khi đó, với mọi $j \in S$:
\[
  \sum_{n=0}^{\infty} P_{i,j}^{(n)} < +\infty
  \qquad \Longrightarrow \qquad
  \lim_{n \to \infty} P_{i,j}^{(n)} = 0.
\]
Lấy tổng theo $j$ trên $S$ hữu hạn:
\[
  \lim_{n \to \infty} \sum_{j \in S} P_{i,j}^{(n)}
  = \sum_{j \in S} \lim_{n \to \infty} P_{i,j}^{(n)} = 0.
\]
Nhưng $\sum_{j \in S} P_{i,j}^{(n)} = 1$ với mọi $n$, nên giới hạn phải bằng $1 \neq 0$.
Mâu thuẫn. Vậy phải tồn tại ít nhất một trạng thái hồi quy.
\end{proof}

% ============================================================================
\section{Hồi quy dương và hồi quy không (Positive and Null Recurrence)}
\label{sec:pos-null-recurrence}

\textbf{Câu chuyện:} Biết rằng bạn chắc chắn quay lại quán café quen --- nhưng MẤT BAO LÂU?
\begin{itemize}[nosep]
  \item \textbf{Hồi quy dương:} Thời gian quay lại trung bình hữu hạn (ví dụ: trung bình 5 ngày quay lại quán).
  \item \textbf{Hồi quy không:} Chắc chắn quay lại, nhưng thời gian trung bình vô hạn (``chờ mãi mãi'').
\end{itemize}

\begin{dinhnghia}[(Thời gian quay lại trung bình -- \en{Mean recurrence time})]
\[
  \mu_i(i) := \E[\tau_i \mid X_0 = i].
\]
\end{dinhnghia}

\begin{dinhnghia}[(Hồi quy dương và hồi quy không)]
Cho trạng thái $i$ là hồi quy ($p_{ii} = 1$).
\begin{itemize}[nosep]
  \item $i$ là \term{hồi quy dương} (\en{positive recurrent}) nếu $\mu_i(i) < +\infty$.
  \item $i$ là \term{hồi quy không} (\en{null recurrent}) nếu $\mu_i(i) = +\infty$.
\end{itemize}
\end{dinhnghia}

\begin{example}[Bước đi ngẫu nhiên đối xứng là hồi quy không]
Bước đi ngẫu nhiên đơn giản trên $\Z$ với $p = q = 1/2$ là \term{hồi quy không}.
Thời gian quay lại trung bình $\mu_0(0) = +\infty$.

Trực quan: bạn \emph{chắc chắn} quay lại gốc, nhưng không ai biết ``trung bình bao lâu'' ---
có thể mất 2 bước, có thể mất 2 triệu bước.
\end{example}

\begin{dinhly}[(Chuỗi hữu hạn và hồi quy dương)]
\label{thm:finite-positive-recurrent}
Trong chuỗi Markov với không gian trạng thái hữu hạn, mọi trạng thái hồi quy đều là
\term{hồi quy dương}.
\end{dinhly}

\begin{corollary}[(Chuỗi bất khả quy hữu hạn)]
\label{cor:finite-irreducible}
Nếu chuỗi Markov là \term{bất khả quy} và có không gian trạng thái hữu hạn,
thì \emph{tất cả} các trạng thái đều là \term{hồi quy dương}.
\end{corollary}

\begin{tomtat}
\begin{center}
\renewcommand{\arraystretch}{1.4}
\begin{tabular}{lccc}
\hline
\textbf{Tính chất} & \textbf{Tạm thời} & \textbf{Hồi quy không} & \textbf{Hồi quy dương} \\
\hline
$p_{ii}$ & $< 1$ & $= 1$ & $= 1$ \\
Quay lại? & Không chắc & Chắc chắn & Chắc chắn \\
$\mu_i(i)$ & --- & $= +\infty$ & $< +\infty$ \\
Ví dụ & RW trên $\Z$, $p \neq q$ & RW trên $\Z$, $p = 1/2$ & Chuỗi hữu hạn \\
\hline
\end{tabular}
\end{center}

\textbf{Quy tắc nhớ:} Chuỗi Markov hữu hạn bất khả quy $\Rightarrow$ mọi trạng thái hồi quy dương.
Hồi quy không chỉ xảy ra khi không gian trạng thái vô hạn.
\end{tomtat}

% ============================================================================
\section{Tính tuần hoàn và phi tuần hoàn (Periodicity and Aperiodicity)}
\label{sec:periodicity}

\textbf{Câu chuyện:} Một số chuỗi Markov ``dao động'' có quy luật. Ví dụ: nếu bạn ở ô trắng trên bàn cờ,
sau 1 bước chắc chắn ở ô đen, sau 2 bước lại ở ô trắng. Chuỗi này có ``nhịp'' (chu kỳ) = 2.
Chuỗi tuần hoàn không hội tụ vì nó cứ ``dao động'' mãi.

\begin{dinhnghia}[(Chu kỳ của trạng thái -- \en{Period of a state})]
\term{Chu kỳ} (\en{period}) của trạng thái $i$ được định nghĩa là:
\begin{equation}\label{eq:period}
  d(i) := \gcd\{n \geq 1 : P_{i,i}^{(n)} > 0\}.
\end{equation}
\end{dinhnghia}

\begin{dinhnghia}[(Phi tuần hoàn -- \en{Aperiodic})]
Trạng thái $i$ được gọi là \term{phi tuần hoàn} (\en{aperiodic}) nếu $d(i) = 1$.

Một chuỗi Markov được gọi là \term{phi tuần hoàn} nếu \emph{tất cả} các trạng thái đều phi tuần hoàn.
\end{dinhnghia}

\textbf{Nhận xét:}
\begin{itemize}[nosep]
  \item Nếu $P_{i,i} > 0$ thì $d(i) = 1$, tức trạng thái $i$ phi tuần hoàn.
        (Mẹo: thêm ``vòng lặp'' tại $i$ là phá vỡ tuần hoàn!)
  \item \term{Trạng thái hấp thụ} ($P_{k,k} = 1$) là đồng thời phi tuần hoàn và hồi quy.
  \item Chu kỳ là tính chất lớp: nếu $i \leftrightarrow j$ thì $d(i) = d(j)$.
\end{itemize}

\begin{example}[Chuỗi có chu kỳ $2$]
Chuỗi Markov trên $S = \{0, 1\}$ với $P = \begin{bmatrix} 0 & 1 \\ 1 & 0 \end{bmatrix}$
có $d(0) = d(1) = 2$: chuỗi ``dao động'' giữa $0$ và $1$.

$P_{0,0}^{(n)} > 0$ chỉ khi $n$ chẵn, nên $\gcd\{2, 4, 6, \ldots\} = 2$.
\end{example}

\begin{example}[Chuỗi có chu kỳ $3$]
Trên $S = \{0, 1, 2\}$ với $P = \begin{bmatrix} 0 & 1 & 0 \\ 0 & 0 & 1 \\ 1 & 0 & 0 \end{bmatrix}$.
Chuỗi quay vòng: $0 \to 1 \to 2 \to 0 \to \cdots$, chu kỳ $d = 3$.
\end{example}

% --- Ergodic definition ---

\begin{dinhnghia}[(Trạng thái ergodic -- \en{Ergodic state})]
Trạng thái $i$ được gọi là \term{ergodic} nếu nó đồng thời:
\begin{enumerate}[label=(\roman*)]
  \item \term{Hồi quy dương}: $p_{ii} = 1$ và $\mu_i(i) < +\infty$.
  \item \term{Phi tuần hoàn}: $d(i) = 1$.
\end{enumerate}
\end{dinhnghia}

\begin{tomtat}
Tính ergodic là điều kiện then chốt để chuỗi Markov hội tụ đến phân phối dừng (Chương~\ref{ch:limiting}).
\[
  \boxed{\text{Ergodic} = \text{Hồi quy dương} + \text{Phi tuần hoàn}}
\]
Chuỗi Markov bất khả quy, hữu hạn và phi tuần hoàn luôn là ergodic.

\medskip
\textbf{Checklist để kiểm tra hội tụ:}
\begin{enumerate}[nosep]
  \item Bất khả quy? (Mọi trạng thái liên thông?)
  \item Phi tuần hoàn? (Có $P_{i,i} > 0$ cho trạng thái nào đó?)
  \item Hữu hạn? (Nếu có $\Rightarrow$ tự động hồi quy dương.)
\end{enumerate}
Nếu cả 3: phân phối dừng tồn tại, duy nhất, và bằng phân phối giới hạn.
\end{tomtat}

% ============================================================================
\section{Bài tập (Exercises)}
\label{sec:ch6-exercises}

\begin{exercise}
\label{ex:ch6-1}
Xét chuỗi Markov trên $S = \{0, 1, 2, 3\}$ với ma trận chuyển:
\[
  P = \begin{bmatrix}
    1/3 & 1/3 & 1/3 & 0 \\
    0   & 0   & 0   & 1 \\
    0   & 1   & 0   & 0 \\
    0   & 0   & 1   & 0
  \end{bmatrix}.
\]
\begin{enumerate}[label=(\alph*)]
  \item Tìm các lớp liên thông.
  \item Tìm chu kỳ của mỗi trạng thái.
  \item Xác định trạng thái nào hồi quy, tạm thời, hấp thụ.
  \item Chuỗi có bất khả quy không?
\end{enumerate}
\end{exercise}

\begin{exercise}
\label{ex:ch6-2}
Chứng minh rằng chu kỳ là tính chất lớp: nếu $i \leftrightarrow j$ thì $d(i) = d(j)$.

(Gợi ý: nếu $P_{i,j}^{(a)} > 0$ và $P_{j,i}^{(b)} > 0$, thì $P_{i,i}^{(a+b)} > 0$,
nên $d(i)$ chia hết $a + b$. Dùng điều này để chỉ $d(i) \mid d(j)$ và ngược lại.)
\end{exercise}

\begin{exercise}
\label{ex:ch6-3}
Xét chuỗi Markov trên $S = \{0, 1, 2, 3, 4\}$ với ma trận chuyển:
\[
  P = \begin{bmatrix}
    0 & 1 & 0 & 0 & 0 \\
    0 & 0 & 1 & 0 & 0 \\
    1/2 & 0 & 0 & 1/2 & 0 \\
    0 & 0 & 0 & 0 & 1 \\
    0 & 0 & 0 & 1 & 0
  \end{bmatrix}.
\]
\begin{enumerate}[label=(\alph*)]
  \item Tìm các lớp liên thông và xác định lớp nào đóng.
  \item Phân loại mỗi trạng thái (hồi quy/tạm thời, tuần hoàn/phi tuần hoàn).
  \item Từ trạng thái 0, chuỗi sẽ kết thúc ở lớp nào?
\end{enumerate}
\end{exercise}

\begin{exercise}
\label{ex:ch6-4}
Cho chuỗi Markov bất khả quy trên không gian trạng thái hữu hạn $S$ với $|S| = N$.
\begin{enumerate}[label=(\alph*)]
  \item Chứng minh rằng mọi trạng thái đều hồi quy.
  \item Chứng minh rằng mọi trạng thái đều hồi quy dương. (Gợi ý: dùng Định lý~\ref{thm:finite-positive-recurrent}.)
  \item Nếu thêm điều kiện phi tuần hoàn, kết luận gì về sự tồn tại phân phối giới hạn?
\end{enumerate}
\end{exercise}
